\documentclass[10pt,a4paper]{amsart}
\usepackage[margin=19mm]{geometry}
\usepackage{amsmath,amssymb,mathtools}
\usepackage[scr=rsfs,cal=euler]{mathalfa}
\usepackage{xfrac}
\usepackage[unicode,hidelinks,bookmarks=false]{hyperref}
\newcommand{\ie}{i.e.\ }
\newcommand{\cf}{cf.\ }
\newcommand{\resp}{respectively\ }
\newcommand{\Subsection}{\subsection}
\providecommand{\nobreakdash}{\nobreak\hbox{-}}
\setlength{\emergencystretch}{2em}
\multlinegap0pt
\usepackage{algorithm}\usepackage{algpseudocode}
\usepackage{amssymb}
\usepackage{mathtools}
\usepackage{xcolor}
\usepackage{tikz}
\usepackage[shortlabels]{enumitem}
\newtheorem{theorem}{Theorem}
\title{Breaking the Orthogonality Barrier in Quantum LDPC Codes}
\author{Kenta Kasai}
\date{Source: arXiv:2601.08824v4 (CC BY 4.0)}
\begin{document}
\maketitle

\noindent\emph{Source and licence.} Breaking the Orthogonality Barrier in Quantum LDPC Codes. Version arXiv:2601.08824v4 (CC BY 4.0), distributed by arXiv under the Creative Commons Attribution 4.0 International licence, http://creativecommons.org/licenses/by/4.0/, as recorded in the arXiv licence metadata for this exact version and shown on the abstract page https://arxiv.org/abs/2601.08824v4. Full licence notice: \texttt{licenses/arxiv-2601.08824-CC-BY-4.0.txt}.\par\smallskip

\noindent\textit{Editorial scope.} Counted unit: the article's theorem theorem (source0.tex lines 234--236). The article's complete original proof of it is delivered with it (source0.tex lines 237--248, one \texttt{proof} environment of 12 lines). No mathematics is written, repaired or re-derived: the statement and the proof are the article's own lines, in the article's own notation, and no retained line is substituted or re-flowed.  No further source-local context is needed: every cross-reference of every retained line resolves inside the two delivered spans.  Nothing was omitted from any retained span, so the package carries no omission record.  No retained line contains a citation, so the wrapper carries no bibliography and no bibliography entry.  No retained line contains a figure environment, an image-inclusion command or drawing code, so no image is delivered and the package declares no asset dependency.  Editorial formatting only, in addition: an AMS \texttt{amsart} preamble that restates the article's own declarations and macros used by the retained lines, namely the environments \texttt{theorem} on separate counters, together with the standard packages those lines need.  The retained environments print in this document as Theorem 1 in order of appearance, whereas the article prints the same items with the numbering of its own full document.  The complete native source of the article as distributed in the arXiv e-print for this exact version is bundled unchanged as \texttt{sources/192632/source0.tex}.
\begin{theorem}
Let $L$ be even and assume $3\le J\le L/2$ with the standard active choice (top $J$ block rows). If $F_u$ and $G_{r-u}$ commute for all $r\in\Delta$ and all $u\in[L/2]$, then the Tanner graphs of the active matrices $H_X$ and $H_Z$ each contain an 8-cycle.
\end{theorem}

\begin{proof}
Since $J\ge 3$, the active block rows include $0,1,2$. Choose $i,j\in[L/2]$ so that $i+j,i+j-1,i+j-2\in\Delta$ (e.g., $i=j=0$). Consider the block positions on the active rows
\[
[(0,i),(0,L/2+j),(1,L/2+j),(1,i+1),(2,i+1),(2,L/2+j+1),(1,L/2+j+1),(1,i)],
\]
where indices are modulo $L/2$. The corresponding cycle product, i.e., the product of permutation matrices encountered along this closed walk, is
\[
W=F_iG_j^{-1}G_{j-1}F_i^{-1}F_{i-1}G_{j-1}^{-1}G_jF_{i-1}^{-1}
\label{eq:8cycle-word-general}
\]
and the pairs $(i,j)$, $(i,j-1)$, $(i-1,j)$, and $(i-1,j-1)$ correspond to the differences $i+j$, $i+j-1$, and $i+j-2$ in $\Delta$. By assumption the relevant $F$ and $G$ commute, so $W=I$ and an 8-cycle exists in the active Tanner graphs.
\end{proof}

\end{document}
